\(\R^2\) 中以 \((0,0)\) 为心、半径为 \(1\) 的开球 \(B_1(0,0)=\{(x,y):x^2+y^2<1\}\)。开集由开球的并生成——这是从距离出发构造拓扑的第一条规则。

本单元依次回答三个问题：

\begin{enumerate}
\tightlist
\item
  距离怎样产生开集与闭集：从度量球定义开、闭、内部、闭包和连续性。
\item
  紧致性把无限局部信息压成有限控制：连接有限子覆盖、收敛子列、极值定理和一致连续。
\item
  完备性让逼近留在空间内：把 Cauchy 判据推广到函数空间，并用 Banach 不动点定理展示完备性怎样使迭代算法有极限。
\end{enumerate}

在 \(\R^n\) 中``闭且有界''等价于紧致，这是有限维欧氏空间的特殊定理，不是紧致性的定义，也不能直接搬到无限维函数空间。

\subsection{一个具体算例}

开覆盖 \([0,1]\subset\bigcup_{x\in[0,1]}(x-0.1,x+0.1)\) 有无限个区间，但只需有限个（例如以 \(0,0.1,0.2,\ldots,1\) 为中心的 11 个区间）就覆盖了 \([0,1]\)——这就是紧致性的有限子覆盖。

\(\Q\) 不是完备的：取有理数列 \(x_1=1,x_2=1.4,x_3=1.41,x_4=1.414,\ldots\) 逐次逼近 \(\sqrt{2}\)，它是 Cauchy 列但在 \(\Q\) 中没有极限。

\subsection{怎么读这一章}

按``邻近---有限控制---极限不逃逸''阅读。距离先产生开集、闭集与连续性，紧致性把无限覆盖压缩成有限子覆盖，完备性则使内部一致的逼近仍落在空间中。紧致与完备解决的不是同一个问题。

\subsection{本章篇目}

以下篇目按概念依赖排列；若从具体问题进入，也可以先打开最接近当前困惑的一篇，再沿篇内链接回补。

\begin{itemize}
\tightlist
\item \hyperref[sec:12-Real-Analysis-Support/03-Topology-Basics/01]{``距离怎样产生开集与闭集''}；
\item \hyperref[sec:12-Real-Analysis-Support/03-Topology-Basics/02]{``紧致性把无限局部信息压成有限控制''}；
\item \hyperref[sec:12-Real-Analysis-Support/03-Topology-Basics/03]{``完备性让逼近留在空间内''}。
\end{itemize}
